\section{Conclusion}
\label{sec:conclusion}

A semantic cache bets that two prompts which look alike can share an answer. We measured how
often that bet loses when no one is attacking it. Separating cache-equivalence from
similarity and labeling it with human judgments, we mapped the false-hit rate of six
embedding models across three domains, traced a false hit into the task it corrupts, and
calibrated the threshold to the cost of being wrong. The picture is sober. On clean domains
a cache can be kept honest, but only by surrendering most of its coverage; on adversarial
near-duplicates the similarity signal carries almost no information about equivalence, and no
threshold makes the cache both safe and useful. The remedy is not a better single number but
a deliberate one: calibrate per domain against a false-hit ceiling set by what a wrong answer
costs, and verify rather than trust. We release the labeled data, the harness, and the exact
commands so that any team can run the same measurement on its own encoder and its own
traffic before it decides what its cache is allowed to guess.
